\section{Black-Scholes Equation}

\subsection{What is an Option?}

\begin{frame}{What is an Option?}
    \begin{itemize}[<+-| alert@+>] % Besides alert, use \pause also works
        \item Derivatives derive value from underlying asset  
        \item Futures, Forwards, Swaps, Options
        \item Options:
        \begin{itemize}
            \item Right, not the obligation, to buy/sell asset at specified price on specified date 
            \item Call gives right to buy
            \item Put gives right to sell
            \item American vs European style
        \end{itemize}
    \end{itemize}
\end{frame}

\subsection{Black-Scholes}

\begin{frame}{Black-Scholes Background}
    \begin{itemize}
        \item Developed in 1970s by Fischer Black, Myron Scholes, and Robert Merton 
        \item Analytic solution to a second order parabolic differential equation
        \item Transformation to heat equation
        \item Options pricing dynamics similar to heat dissipation
        \item Based on European style options
        \item American style options pricing can only be done numerically, not analytically
    \end{itemize}
\end{frame}

\begin{frame}{Black-Scholes Equation}
    \begin{itemize}
        \item 
    \end{itemize}
\end{frame}

\subsection{Boundary Conditions}

\begin{frame}{Boundary Conditions}
    \begin{itemize}
        \item 
        \item 
        \item 
    \end{itemize}
\end{frame}

\subsection{Transformation to the Heat Equations}

\begin{frame}{Transformation to the Heat Equations}
    \begin{itemize}
        \item 
        \item 
        \item 
    \end{itemize}
\end{frame}
